\section{Introduction}

Multichannel surface-wave methods use spatially distributed vibration measurements to estimate the dispersive behavior of Rayleigh waves and infer near-surface mechanical properties \cite{Park1999,Foti2014}. Although the processing methods are well established, the acquisition hardware remains a significant part of the cost and deployment complexity, so compact wireless nodes built from inexpensive geophones and mixed-signal microcontrollers are attractive for dense, reconfigurable arrays \cite{Kafadar2020}.

The present work forms part of a low-cost Internet-of-Things acquisition platform based on SM-24 moving-coil geophones \cite{SM24}. A PSoC~5LP was selected as the analog front end because it integrates programmable-gain amplifiers, operational amplifiers, analog filters, voltage DACs, analog multiplexers, a delta-sigma ADC, a digital filter block, DMA, and an Arm Cortex-M3 processor in one device \cite{InfineonPSoC}. The guiding principle is to reuse this already-available hardware rather than add dedicated calibration components: the same ADC and digital filter are time-shared between acquisition and calibration, so the correction loop uses only resources the platform already carries, avoiding a separate converter or processor. The same reconfigurability also imposes a limitation: its internal analog blocks cannot be independently optimized for low offset, drift, and flicker noise -- Lopin and Lopin reported the same tradeoff in a single-chip PSoC potentiostat, with poorer offset and noise specifications than purpose-selected external amplifiers \cite{Lopin2018}.

The geophone conditioning circuit is derived from the cascade correction and subtraction topology of Ma \emph{et al.} \cite{Ma2023}; mapped to the PSoC, accumulated DC errors prevented sufficiently aggressive gain settings and became particularly severe at the ADDER output.

A continuously active DC servo, evaluated as an initial solution, proved impractical (Section~\ref{sec:discussion}); the proposed method instead performs a finite foreground calibration before measurement and disconnects the controller afterward.

Self-calibration using on-chip resources is not itself new. The PSoC-Stat injects five known current levels with an IDAC and uses the resulting ADC values to estimate the gain and offset of a TIA/ADC signal chain \cite{Lopin2018}. Infineon's AN68403 applies an internal VDAC to a PGA/ADC chain, measures the system offset with grounded inputs, and stores gain and offset corrections in EEPROM \cite{InfineonAN68403}. These methods primarily correct the numerical transfer characteristic observed at the ADC output. The objective here is different: the proposed loop changes the physical reference of each amplification stage so that its output is adjusted toward its intended operating point, within the available DAC range and resolution.

The front end is deliberately DC-coupled rather than AC-coupled: the SM-24's usable band extends down toward its mechanical corner, and a coupling capacitor placed ahead of it would need a corner well below 1~Hz to avoid attenuating exactly the low-frequency content the acquisition targets, degrading the signal it is meant to preserve. That constraint rules out the simplest fix for offset -- block it with a capacitor -- and leaves per-stage trimming as the remaining option once the amplifiers themselves cannot be swapped for lower-offset external parts without abandoning the single-chip premise.

An ADDER-only fixed DC injection was tried first, as an early workaround (Section~\ref{sec:platform-design}): it recenters one node but cannot recover stages before or after it.

The contribution is threefold: 1) a per-stage analog calibration architecture that reuses the PSoC analog multiplexer, ADC, digital filter, and four VDACs, so that improving an existing budget-constrained analog front end does not require oversizing each stage with dedicated precision components; 2) a measurement-gated foreground implementation, demonstrated with a gain-aware PI controller using deadband, anti-windup, slew limiting, stability detection, and one-LSB refinement, that freezes all corrections before seismic data are recorded and so introduces no additional closed-loop poles into the acquisition transfer function; and 3) an experimental validation on two independently assembled boards across the PGA's full configurable gain range ($\times1$ to $\times50$), showing that the calibrated front end keeps every stage within a small fraction of the ADC's full-scale range regardless of board or gain, even where an individual stage such as the PGA converges only partially. The design goal is not to null every stage's offset exactly, but to guarantee that no board, at no gain setting, requires expert per-unit tuning to stay within its acceptance range.

